%======================================================================
% فصل پنجم: ارزیابی
%======================================================================
\chapter{ارزیابی}
\label{ch:evaluation}

\section{مقدمه}

هدف از ارزیابی این پروژه، سنجش میزان انطباق پیاده‌سازی با اهداف تعریف‌شده
در فصل اول است. از آن‌جا که سامانه‌ی حاضر یک محصول زیرساختی است، ارزیابی
آن صرفاً بر مبنای کیفیت پاسخ مدل‌ها انجام نمی‌شود؛ بلکه معیارهایی مانند
درستی جریان‌های مالی، استحکام لایه‌ی دسترسی، پوشش آزمون، قابلیت توسعه و
ویژگی‌های عملیاتی نیز اهمیت دارند.

در این فصل، ارزیابی در سه محور انجام می‌شود: ارزیابی فنی مبتنی بر آزمون‌ها،
ارزیابی عملکرد معماری، و مقایسه‌ی دستاوردهای نهایی با اهداف اولیه.

\section{روش ارزیابی}

ارزیابی این پایان‌نامه بر سه منبع شواهد استوار است. منبع نخست، اجرای واقعی
مجموعه‌آزمون خودکار برای سنجش درستی رفتار سرویس‌ها، نقاط پایانی و
سناریوهای انتهابه‌انتها است. منبع دوم، بازبینی قراردادهای قابل مشاهده‌ی
سامانه است تا روشن شود دامنه‌ی \lr{API} مستندشده با پیاده‌سازی واقعی
سازگار است. منبع سوم نیز تحلیل مسیر بحرانی معماری است که نشان می‌دهد کدام
بخش‌ها در مسیر هم‌زمان درخواست باقی مانده‌اند و کدام بخش‌ها به کش و صف
منتقل شده‌اند.

این رویکرد آگاهانه انتخاب شده است؛ زیرا موضوع پایان‌نامه بیشتر
\textbf{درستی نرم‌افزاری و کفایت معماری} را هدف می‌گیرد تا مقایسه‌ی
سرعت مدل‌های هوش مصنوعی. در نتیجه، معیارهای این فصل به‌جای کیفیت پاسخ
مدل‌ها، بر صحت قرارداد، پایداری مسیرهای مالی و قابلیت اتکای سامانه
تمرکز دارند.

\subsection{محیط اجرا و داده‌های پایه}

اجرای اصلی ارزیابی در تاریخ 2026-03-23 با فرمان \code{php artisan test}
در محیط توسعه‌ی محلی پروژه انجام شده است. جدول
\ref{tab:evaluation-setup} داده‌های پایه‌ی این اجرا را نشان می‌دهد.

\begin{table}[H]
\centering
\caption{مشخصات اجرای ارزیابی}
\label{tab:evaluation-setup}
\renewcommand{\arraystretch}{1.4}
\begin{tabularx}{\textwidth}{>{\raggedleft\arraybackslash}p{4.2cm} X}
\toprule
\textbf{شاخص} & \textbf{مقدار} \\
\midrule
نسخه‌ی بستر سمت سرور & \lr{PHP 8.2} و \lr{Laravel 12} \\
فرمان ارزیابی اصلی & \code{php artisan test} \\
نتیجه‌ی کلی اجرا & موفق، بدون شکست در مجموعه‌آزمون \\
تعداد کل آزمون‌ها & 55 آزمون \\
تعداد کل \lr{Assertion}ها & 160 \\
زمان اجرای مجموعه‌آزمون & 65.91 ثانیه \\
دامنه‌ی \lr{API} نسخه‌بندی‌شده & 35 مسیر شامل 6 مسیر هوش مصنوعی و 29 مسیر مدیریتی، مالی و گزارشی \\
\bottomrule
\end{tabularx}
\end{table}

\section{ارزیابی فنی}

\subsection{آزمون‌های واحد}

برای سنجش درستی اجزای مستقل سامانه، مجموعه‌ای از آزمون‌های واحد تعریف
شده است. این آزمون‌ها روی سرویس‌ها، کارهای صف، سیاست‌های کش و اجزای
ثبت گزارش متمرکز هستند. در کدپایه‌ی فعلی، ۱۳ فایل آزمون واحد وجود دارد
که بخش‌های اصلی مختلفی را پوشش می‌دهند. برای مثال،
\code{WalletService} و \code{SubscriptionService} برای اطمینان از درستی
منطق مالی آزموده شده‌اند، \code{InvoiceService} ساخت فاکتور و اقلام آن را
پوشش می‌دهد، \code{UsageEstimationService} و
\code{UsageMeteringService} به تخمین و محاسبه‌ی هزینه اختصاص دارند،
\code{GatewayLogService} و مقصدهای گزارش خارجی پایداری لاگینگ را می‌سنجند
و کارهای پس‌زمینه‌ای مانند تجمیع مصرف روزانه و پایش سلامت ارائه‌دهنده نیز
به‌طور مستقل بررسی شده‌اند.

تمرکز این دسته از آزمون‌ها بر درستی منطق و رفتار اجزای منفرد است؛ بنابراین
می‌توان خطاهای مربوط به محاسبه، تبدیل داده یا سیاست‌های کنترلی را پیش از
درگیر شدن کل سامانه تشخیص داد.

\subsection{آزمون‌های ویژگی}

در کنار آزمون‌های واحد، ۱۵ فایل آزمون ویژگی برای سنجش رفتار
\lr{API} و سناریوهای انتهابه‌انتها تعریف شده است. این آزمون‌ها شامل
گروه‌های مهمی مانند احراز هویت \lr{OTP}، پروفایل و کش پروفایل، مدیریت
مشتری و کلید \lr{API}، کیف پول، طرح‌ها و اشتراک‌ها، فاکتورها و گزارش‌ها،
مدیریت ارائه‌دهنده و مدل، نقاط پایانی اصلی درگاه هوش مصنوعی و همچنین لاگ
ممیزی و سناریوهای کاربری ترکیبی را در بر می‌گیرند.

وجود این آزمون‌ها نشان می‌دهد که طراحی پروژه فقط در سطح کلاس‌ها متوقف
نشده و رفتار واقعی سامانه در مسیرهای مهم نیز مورد بررسی قرار گرفته است.

\subsection{نتایج آزمون}

به‌جای اکتفا به توصیف کیفی، نتیجه‌ی اجرای واقعی مجموعه‌آزمون در
جدول \ref{tab:test-summary} گزارش شده است. این داده‌ها نشان می‌دهند که
پیاده‌سازی فعلی، دست‌کم در محدوده‌ی سناریوهای تعریف‌شده، از پایداری
نرم‌افزاری مناسبی برخوردار است.

\begin{table}[H]
\centering
\caption{خلاصه‌ی اجرای مجموعه‌آزمون}
\label{tab:test-summary}
\renewcommand{\arraystretch}{1.4}
\begin{tabularx}{\textwidth}{>{\raggedleft\arraybackslash}p{4cm} >{\centering\arraybackslash}p{2.5cm} X}
\toprule
\textbf{شاخص} & \textbf{مقدار} & \textbf{توضیح} \\
\midrule
فایل‌های آزمون واحد & ۱۳ & پوشش سرویس‌های مالی، تخمین مصرف، کش، لاگینگ و کارهای پس‌زمینه \\
فایل‌های آزمون ویژگی & ۱۵ & پوشش احراز هویت، مدیریت کلید، کیف پول، اشتراک، گزارش و مسیرهای درگاه \\
تعداد کل آزمون‌ها & 55 & شامل آزمون‌های واحد و ویژگی \\
تعداد کل \lr{Assertion}ها & 160 & نشان‌دهنده‌ی آن است که هر آزمون به یک بررسی سطحی محدود نشده است \\
وضعیت نهایی & موفق & همه‌ی آزمون‌ها در اجرای داوری با موفقیت عبور کرده‌اند \\
\bottomrule
\end{tabularx}
\end{table}

\subsection{سناریوهای کلیدی تأییدشده}

برای ارزیابی یک سامانه‌ی واسط، صرف تعداد آزمون‌ها کافی نیست؛ مهم‌تر آن است
که آیا مسیرهای بحرانی واقعاً سنجیده شده‌اند یا خیر. جدول
\ref{tab:key-scenarios} تعدادی از سناریوهای مهمی را که در کدپایه به‌صورت
خودکار بررسی شده‌اند نشان می‌دهد.

\begin{table}[H]
\centering
\caption{نمونه‌ای از سناریوهای کلیدی تأییدشده در آزمون‌ها}
\label{tab:key-scenarios}
\renewcommand{\arraystretch}{1.4}
\begin{tabularx}{\textwidth}{>{\raggedleft\arraybackslash}p{4.2cm} >{\raggedleft\arraybackslash}p{3.2cm} X}
\toprule
\textbf{سناریو} & \textbf{شاهد آزمون} & \textbf{نتیجه} \\
\midrule
شروع و تأیید احراز هویت \lr{OTP} & \code{AuthOtpTest} & جریان تولید چالش، تأیید کد و صدور توکن نشست بررسی شده است. \\
رد درخواست فاقد کلید یا دارای داده‌ی نامعتبر & \code{GatewayTest} & درگاه برای نبود کلید، نبود ارائه‌دهنده و خطای اعتبارسنجی پاسخ استاندارد تولید می‌کند. \\
جلوگیری از مصرف بدون موجودی کافی & \code{UserScenarioTest} & در صورت کمبود موجودی، درخواست پیش از فراخوانی ارائه‌دهنده متوقف می‌شود. \\
پشتیبانی از مصرف مبتنی بر اشتراک & \code{UserScenarioTest} & کاربر دارای اشتراک معتبر می‌تواند بدون برداشت مستقیم از کیف پول از سرویس استفاده کند. \\
نرمال‌سازی تفاوت ارائه‌دهندگان & \code{GatewayTest} & پاسخ‌های \lr{Gemini} به قالب سازگار با \lr{OpenAI} برگردانده می‌شوند و مسیر \lr{Groq} نیز بررسی شده است. \\
تولید خروجی‌های جانبی و قابلیت استفاده‌ی ابزاری & \code{InvoiceTest} و \code{UserScenarioTest} & تولید \lr{PDF} فاکتور و در دسترس بودن \lr{Playground} تأیید شده‌اند. \\
\bottomrule
\end{tabularx}
\end{table}

\section{ارزیابی عملکرد}

\subsection{تأخیر پاسخ}

در این پایان‌نامه بنچ‌مارک میدانی مستقل برای اندازه‌گیری تأخیر و توان عملیاتی
انجام نشده است؛ بنابراین ارائه‌ی اعداد مطلق برای \lr{latency} یا
\lr{throughput} از نظر روش‌شناختی موجه نیست. با این حال، تحلیل مسیر بحرانی
سامانه نشان می‌دهد که طراحی انجام‌شده تلاش کرده است سربار درون‌سامانه‌ای را
در بخش‌های قابل‌کنترل کاهش دهد. برای مثال، داده‌های پرتکرار مانند
ارائه‌دهندگان، مدل‌ها، طرح‌ها و پروفایل کاربر از کش خوانده می‌شوند،
کارهای غیر بحرانی مانند ارسال گزارش خارجی، تولید \lr{PDF} و ارسال
\lr{OTP} از مسیر اصلی درخواست خارج و به صف منتقل شده‌اند، و استفاده از
الگوی خط لوله سربار منطقی سامانه را به مراحل کوچک و ساده تقسیم کرده و
پیچیدگی هر مرحله را محدود نگه داشته است.

بر این اساس، می‌توان گفت که معماری سامانه برای کاهش سربار مسیر اصلی
\textbf{طراحی شده است}، اما ادعای عددی درباره‌ی کارایی نهایی آن نیازمند
آزمون بار مستقل در محیطی نزدیک به بهره‌برداری است. این تمایز برای پرهیز
از بیش‌اظهاری در ارزیابی ضروری است.

\subsection{دقت صورت‌حساب‌گذاری}

دقت مالی در این پروژه از دو مسیر تقویت شده است. سامانه پیش از ارسال
درخواست، با استفاده از توکن‌ساز \lr{BPE} و معیارهای غیرتوکنی، هزینه‌ی
احتمالی را محاسبه می‌کند تا بتواند تصمیم مالی اولیه را بگیرد. پس از
دریافت پاسخ نیز، اگر داده‌ی مصرف از سوی ارائه‌دهنده موجود باشد همان مقدار
واقعی ملاک قرار می‌گیرد و در غیر این صورت، تخمین اولیه به‌عنوان پشتیبان
باقی می‌ماند.

این طراحی دو مزیت مهم دارد. نخست، از ارسال درخواست‌هایی که موجودی کافی
ندارند جلوگیری می‌شود. دوم، حتی در صورت نبود فراساختار مصرف در پاسخ
ارائه‌دهنده، سامانه همچنان قادر است مدل مالی یکنواختی ارائه دهد. از منظر
شواهد نرم‌افزاری، آزمون‌های \code{UsageEstimationServiceTest}،
\code{WalletServiceTest}، \code{SubscriptionServiceTest} و
\code{UserScenarioTest} این منطق را در سطح سرویس و سناریوی کاربری بررسی
می‌کنند. با این حال، تطبیق نهایی با صورت‌حساب واقعی همه‌ی ارائه‌دهندگان
هنوز نیازمند اعتبارسنجی میدانی و داده‌های عملیاتی بلندمدت است.

\section{مقایسه با اهداف}

در فصل اول، شش هدف عملیاتی برای پروژه تعریف شد. جدول
\ref{tab:goal-comparison} میزان تحقق این اهداف را نشان می‌دهد.

\begin{table}[H]
\centering
\caption{مقایسه‌ی اهداف اولیه با وضعیت نهایی پیاده‌سازی}
\label{tab:goal-comparison}
\renewcommand{\arraystretch}{1.4}
\begin{tabularx}{\textwidth}{>{\raggedleft\arraybackslash}p{4cm} >{\centering\arraybackslash}p{2cm} X}
\toprule
\textbf{هدف} & \textbf{وضعیت} & \textbf{شاهد و توضیح} \\
\midrule
ارائه‌ی رابط یکپارچه‌ی سازگار با \lr{OpenAI} & کامل & 6 مسیر اصلی هوش مصنوعی در \lr{API} موجود است و آزمون‌های \code{GatewayTest} رفتار آن‌ها را بررسی می‌کنند. \\
پشتیبانی از چند ارائه‌دهنده & کامل & چهار ارائه‌دهنده پشتیبانی می‌شوند؛ با این توضیح که پوشش قابلیت‌های \lr{Gemini} در نسخه‌ی فعلی به گفتگو، \lr{responses} و جاسازی متن محدود است. \\
احراز هویت و مدیریت دسترسی & کامل & مسیرهای \lr{OTP}، مشتریان \lr{API}، کلیدها و لاگ ممیزی وجود دارند و در آزمون‌های ویژگی پوشش داده شده‌اند. \\
سیستم صورت‌حساب و کیف پول & کامل & سرویس‌های کیف پول، اشتراک و فاکتور به‌همراه سناریوهای بدهکارسازی و کنترل موجودی آزموده شده‌اند. \\
گزارش‌دهی و مانیتورینگ & کامل & گزارش مصرف، دفترکل کیف پول، گزارش فاکتورها و دو مقصد گزارش خارجی در کدپایه و آزمون‌ها حاضر هستند. \\
زیرساخت عملیاتی کامل بیرونی & نسبی & منطق ارسال \lr{OTP} و تولید \lr{PDF} وجود دارد، اما اتصال به سرویس واقعی پیامک/ایمیل و درگاه پرداخت بیرونی هنوز تکمیل نشده است. \\
\bottomrule
\end{tabularx}
\end{table}

بر پایه‌ی این مقایسه می‌توان گفت که اهداف اصلی پروژه محقق شده‌اند و
نقاط باقی‌مانده بیشتر به یکپارچه‌سازی با سرویس‌های عملیاتی بیرونی مربوط
می‌شوند، نه به منطق مرکزی سامانه. در عین حال، اعتبار این نتیجه‌گیری
بیشتر در حوزه‌ی \textbf{درستی رفتاری و انسجام معماری} است و نه در حوزه‌ی
\textbf{کارایی میدانی در مقیاس تولید}.
